% Generated draft: requires mathematical and source review.
\documentclass[11pt]{article}
\usepackage[margin=25mm]{geometry}
\usepackage{fontspec}
\setmainfont{Latin Modern Roman}
\usepackage{amsmath,amssymb,mathtools,unicode-math}
\setmathfont{Latin Modern Math}
\usepackage{enumitem,needspace,xurl,hyperref}
\setlength{\parindent}{0pt}
\setlength{\parskip}{5pt}
\setlength{\emergencystretch}{3em}
\newcommand{\N}{\mathbb{N}}
\newcommand{\Z}{\mathbb{Z}}
\newcommand{\Q}{\mathbb{Q}}
\newcommand{\R}{\mathbb{R}}
\newcommand{\C}{\mathbb{C}}
\newcommand{\F}{\mathbb{F}}
\newcommand{\A}{\mathbb{A}}
\newcommand{\PP}{\mathbb{P}}
\newcommand{\E}{\mathbb{E}}
\newcommand{\eps}{\varepsilon}
\newcommand{\dd}{\,\mathrm{d}}
\begin{document}
\section*{UnsolvedMath ID: 30004080}
\textbf{Source problem code:} catalogue entry
\subsection{Definitions and mathematical statement}
Let \(k\) be an \(F\)-finite field of characteristic \(p>0\), meaning that \([k:k^p]<\infty\). Let \(A\) be a normal domain of finite type over \(k\), and suppose that \(A\) is numerically \(\Q\)-Gorenstein. Set \(R=S^{-1}A\), where \(S\subseteq A\setminus\{0\}\) is a multiplicatively closed subset containing \(1\). Here the numerically \(\Q\)-Gorenstein hypothesis means that a canonical Weil divisor is numerically \(\Q\)-Cartier. The supplied record does not specify the precise numerical pullback convention underlying this hypothesis; its intended meaning must be checked against the cited source. It must not be replaced silently by the stronger \(\Q\)-Gorenstein hypothesis.

Fix an algebraic closure \(\overline{K}\) of \(K=\operatorname{Frac}(R)\). The absolute integral closure \(R^+\) is the integral closure of \(R\) in \(\overline{K}\). For each ideal \(I\) of \(R\), define its plus closure and the relevant colon ideal by
\[
I^+=IR^+\cap R,
\qquad
(I:I^+)=\{r\in R:rI^+\subseteq I\}.
\]
Define
\[
P(R)=\bigcap_{I\text{ an ideal of }R}(I:I^+).
\]
The intersection ranges over all ideals of \(R\), with no restriction to parameter ideals or primary ideals.

For clarity, use the standard big tight closure convention for the big test ideal. For any \(R\)-module \(M\), with no finite generation assumption, let
\[
F_R^e(M)=R\otimes_{F^e,R}M,
\]
where the right \(R\)-action on the first factor is \(a\cdot r=ar^{p^e}\). Define
\[
0_M^*=\left\{m\in M:\text{ there exists }c\in R\setminus\{0\}\text{ such that }c(1\otimes m)=0\text{ in }F_R^e(M)\text{ for all sufficiently large }e\right\}.
\]
Then the big test ideal is
\[
\tau_b(R)=\bigcap_{M\text{ an }R\text{-module}}\operatorname{Ann}_R(0_M^*),
\qquad
\operatorname{Ann}_R(N)=\{r\in R:rN=0\}.
\]
\textbf{Question.} For every \(R\) satisfying the stated hypotheses, is
\[
\tau_b(R)=P(R)?
\]
The source labels this question open. That label and the supplied dated literature assessment are source information, not an independently verified mathematical conclusion.
\subsection{Short English statement}
For each ring in the stated numerically \(\Q\)-Gorenstein class, intersect the ideals of elements that send every plus closure back into its original ideal. Must this intersection equal the big test ideal? The precise numerical \(\Q\)-Gorenstein convention remains to be checked in the cited source.
\subsection{Sources}
\begin{itemize}
\item Supplied local registry record: \texttt{data/unsolvedmath/problems.json}, permanent UnsolvedMath ID \(30004080\), title \emph{Big Test Ideals from Plus Closure}, source problem code \texttt{catalogue entry}, in the Oberwolfach Reports Open Problems collection.
\item Supplied report citation: \emph{Singularities and Homological Aspects of Commutative Algebra (2020)}, Oberwolfach Reports, pp.~407--464, DOI \texttt{10.4171/owr/2019/6}. \url{https://doi.org/10.4171/owr/2019/6}. The literature-source label is \emph{Oberwolfach Report 6/2019}; the record gives proposed year 2019 and names no individual proposer.
\item Supplied literature reference, labeled \emph{Takagi, finitistic test ideals}: \url{https://doi.org/10.1016/j.jalgebra.2018.08.003}. The supplied record gives no fuller bibliographic description or precise supporting location.
\item Literature assessment dated exactly \texttt{2026-08-22}: the supplied curation labels the problem open and states that the exact equality remains unproved, while known inclusions and related special-case equalities do not settle it. This is supplied dated curation derived from the OWR material, not an independently verified assessment; time-sensitive status and cited claims require independent verification.
\end{itemize}
\end{document}
